% ============================================================
% CAPITOLUL 5 — IMPLEMENTAREA MODELELOR SIDARTH ȘI SIR
% ============================================================

\chapter{Implementarea modelelor SIDARTH și SIR}
\label{ch:implementare}

\section{Datele de intrare și inițializarea compartimentelor}
\label{sec:seturi_date}

Datele epidemiologice utilizate de sistem provin din setul public
\emph{Our World in Data} (OWID) COVID-19 \cite{owid_covid}. Fișierele brute
nu conțin direct toate compartimentele necesare unui model SIDARTH, ci serii
temporale agregate, cum ar fi cazurile noi, cazurile cumulate, decesele
cumulate și populația. Din acest motiv, etapa de preprocesare nu pretinde că
recuperează starea epidemiologică completă a populației. Ea construiește o
stare inițială compatibilă cu sistemul software, folosind datele observabile
drept punct de pornire pentru explorarea numerică a scenariilor.

Scriptul \texttt{scripts/preprocess\_owid.py} reindexează fiecare țară pe un
calendar zilnic continuu, completează cumulativii prin forward-fill și înlocuiește
valorile negative din \texttt{new\_cases} cu zero, deoarece acestea apar de obicei
ca ajustări retrospective ale raportărilor. Numărul de infectați activi este
aproximat printr-o fereastră mobilă de 14 zile:
\begin{equation}
I(t) = \sum_{\tau=t-\Delta}^{t} \text{new\_cases}(\tau), \qquad \Delta = 14.
\end{equation}
Compartimentul \emph{removed} este aproximat prin
\begin{equation}
R_{\text{removed}}(t)=\text{total\_cases}(t)-I(t),
\end{equation}
iar susceptibilii se obțin din bilanțul populației,
\begin{equation}
S(t)=N-I(t)-R_{\text{removed}}(t).
\end{equation}

Pentru modelul SIR, aceste trei valori pot fi folosite direct ca
$(S_0,I_0,R_0)$. Pentru SIDARTH, aceeași informație este distribuită într-o
stare inițială mai bogată: cazurile active observabile sunt plasate în
compartimentul $I$, persoanele ieșite din infecție sunt plasate în $H$, iar
compartimentele $D$, $A$, $R$ și $T$ pornesc de la zero. Această alegere este
deliberat conservatoare. Ea nu afirmă că distribuția clinică reală este
cunoscută, ci oferă o condiție inițială reproductibilă, potrivită pentru
compararea implementărilor și pentru analiza sensibilității la parametri.

\begin{table}[htbp]
\centering
\caption{Maparea datelor observabile în condiții inițiale pentru modelele suportate.}
\label{tab:initializare_modele}
\begin{tabular}{lll}
\toprule
\textbf{Model} & \textbf{Compartimente inițiale} & \textbf{Interpretare} \\
\midrule
SIR & $S_0$, $I_0$, $R_0$ & model agregat cu stare minimă \\
SIDARTH & $S_0$, $I_0$, $D_0=A_0=R_0=T_0=0$, $H_0$ & model clinic stratificat \\
\bottomrule
\end{tabular}
\end{table}

\section{Contractul comun de simulare}
\label{sec:contract_simulare}

Pivotarea către SIDARTH a impus schimbarea contractului intern dintre API și
runnere. Versiunea inițială a sistemului transmitea doar parametrii
$\beta$ și $\gamma$ și presupunea implicit că rezultatul conține curbele
$S$, $I$ și $R$. Noua versiune introduce câmpul \texttt{model\_type}, un obiect
\texttt{parameters} și un obiect \texttt{initial}. În acest fel, același flux
poate executa SIDARTH ca model principal și SIR ca model secundar, fără a
modifica interfața de cozi sau mecanismul de benchmark.

\begin{lstlisting}[language=bash, caption={Structura job-ului publicat în Redis pentru o simulare SIDARTH.}]
{
  "simulation_id": "a3f9c2d1-...",
  "dataset_id": "7b4e1a02-...",
  "model_type": "sidarth",
  "language": "python",
  "parameters": {
    "alpha": 0.57, "beta": 0.011, "gamma": 0.456, "delta": 0.011,
    "epsilon": 0.171, "zeta": 0.125, "lambda": 0.034,
    "eta": 0.125, "rho": 0.034, "theta": 0.370,
    "mu": 0.017, "kappa": 0.017, "nu": 0.027,
    "xi": 0.017, "sigma": 0.017
  },
  "initial": {
    "s": 19000000, "i": 1200, "d": 0, "a": 0,
    "r": 0, "t": 0, "h": 18000
  },
  "population_n": 19019200,
  "t_max": 160,
  "db_dsn": "postgres://..."
}
\end{lstlisting}

Coada Redis a fost redenumită din \texttt{sir:jobs:<language>} în
\texttt{epimodel:jobs:<language>}. Schimbarea este mică în cod, dar importantă
conceptual: runner-ul nu mai este o implementare exclusivă SIR, ci un executor
pentru modele compartimentale compatibile cu schema comună. Rezultatele sunt
stocate în \texttt{simulation\_results} ca obiect JSONB \texttt{compartments},
nu ca trei coloane fixe. Pentru o rulare SIDARTH, un punct temporal conține
cheile \texttt{s}, \texttt{i}, \texttt{d}, \texttt{a}, \texttt{r}, \texttt{t}
și \texttt{h}; pentru SIR, același câmp conține doar \texttt{s}, \texttt{i}
și \texttt{r}.

\section{Implementarea modelului SIDARTH}
\label{sec:implementare_sidarth}

Modelul implementat în această lucrare este o variantă SIDARTH derivată din
familia SIDARTHE propusă de Giordano et al. \cite{giordano2020sidarthe}. Se
păstrează compartimentele $S$, $I$, $D$, $A$, $R$, $T$ și $H$, iar
compartimentul $E$ al deceselor este omis pentru a concentra aplicația pe
dinamica infecției, diagnosticării, severității și vindecării. Ecuația pentru
$T$ folosește astfel doar ieșirea către $H$, prin rata $\sigma$, nu și termenul
de mortalitate $\tau$.

În Python, nucleul modelului este o funcție RHS compatibilă cu
\texttt{scipy.integrate.solve\_ivp}. Ea primește vectorul de stare, dicționarul
de parametri și populația totală:
\begin{lstlisting}[language=Python, caption={Funcția RHS SIDARTH în runner-ul Python.}]
def sidarth_rhs(t, y, params, n):
    s, i, d, a, r, threatened, h = y
    force = (
        params["alpha"] * i +
        params["beta"] * d +
        params["gamma"] * a +
        params["delta"] * r
    ) / n
    new_infections = s * force

    return [
        -new_infections,
        new_infections - (params["epsilon"] + params["zeta"] + params["lambda"]) * i,
        params["epsilon"] * i - (params["eta"] + params["rho"]) * d,
        params["zeta"] * i - (params["theta"] + params["mu"] + params["kappa"]) * a,
        params["eta"] * d + params["theta"] * a - (params["nu"] + params["xi"]) * r,
        params["mu"] * a + params["nu"] * r - params["sigma"] * threatened,
        params["lambda"] * i + params["rho"] * d + params["kappa"] * a
            + params["xi"] * r + params["sigma"] * threatened,
    ]
\end{lstlisting}

Aceeași structură este implementată în R și Julia. În R, funcția este transmisă
către \texttt{deSolve::ode()}, cu metoda \texttt{rk4}, păstrând pasul zilnic fix.
În Julia, funcția respectă convenția \emph{in-place} din
\texttt{DifferentialEquations.jl}, iar integrarea folosește \texttt{Tsit5()} cu
\texttt{saveat = 0:1.0:t\_max}. Diferența dintre cele trei implementări nu este
matematică, ci ține de ecosistemul numeric: R oferă o interfață stabilă și
compactă pentru ODE-uri clasice, Julia urmărește performanța și controlul fin
al solverului, iar Python oferă integrare directă cu infrastructura software și
cu bibliotecile științifice uzuale.

\section{SIR ca model secundar}
\label{sec:sir_secundar}

Modelul SIR rămâne în sistem pentru comparații și pentru verificări de regresie.
El are un rol diferit față de versiunea inițială a lucrării. Nu mai este modelul
principal prin care se discută dinamica clinică a epidemiei, ci un reper
minimal: dacă două limbaje produc rezultate apropiate pentru SIR, diferențele
observate ulterior în SIDARTH pot fi analizate în raport cu complexitatea
modelului, nu cu erori elementare de integrare sau de persistare.

În contractul comun, SIR folosește același câmp \texttt{parameters}, dar cu
doar două rate, \texttt{beta} și \texttt{gamma}. Condițiile inițiale sunt
transmise prin \texttt{initial = \{s, i, r\}}, iar rezultatele sunt stocate tot
în câmpul \texttt{compartments}. Această simetrie simplifică interfața web:
graficul citește cheile disponibile în fiecare punct temporal și construiește
curbele fără să presupună apriori numărul de compartimente.

\section{Persistența rezultatelor și exportul}
\label{sec:persistenta_rezultate}

Tabela \texttt{simulation\_runs} reține metadatele execuției: utilizatorul,
dataset-ul, modelul, limbajul, populația, orizontul temporal și parametrii sub
formă JSONB. Tabela \texttt{simulation\_results} reține câte un rând pentru
fiecare zi simulată, identificat prin perechea \texttt{(simulation\_id, t)}.
În locul coloanelor fixe \texttt{s}, \texttt{i}, \texttt{r}, rândul conține
obiectul JSONB \texttt{compartments}. Această decizie evită migrații repetate
ale schemei atunci când se adaugă un model cu altă structură de stare.

\begin{lstlisting}[language=SQL, caption={Schema generică pentru rezultatele simulărilor.}]
CREATE TABLE simulation_results (
    simulation_id UUID NOT NULL REFERENCES simulation_runs(id) ON DELETE CASCADE,
    t             INTEGER NOT NULL,
    compartments  JSONB NOT NULL,
    PRIMARY KEY (simulation_id, t),
    CHECK (t >= 0),
    CHECK (jsonb_typeof(compartments) = 'object')
);
\end{lstlisting}

Exportul CSV este generat dinamic. API-ul inspectează cheile prezente în
rezultate, le ordonează în secvența epidemiologică $S,I,D,A,R,T,H$ și scrie
doar coloanele existente pentru modelul respectiv. Prin urmare, o simulare SIR
se exportă cu coloanele \texttt{t,s,i,r}, iar o simulare SIDARTH cu
\texttt{t,s,i,d,a,r,t,h}. Această convenție păstrează fișierele lizibile și
evită completarea artificială a unor compartimente care nu există în model.

\section{Metrici de performanță}
\label{sec:metrici_implementare}

Toate runnerele colectează aceleași trei metrici: timpul wall-clock, timpul CPU
și vârful de memorie rezidentă. Python folosește \texttt{time.perf\_counter()},
\texttt{time.process\_time()} și \texttt{resource.getrusage}; R folosește
\texttt{proc.time()} și \texttt{/proc/self/status}; Julia folosește
\texttt{time\_ns()} și \texttt{Sys.maxrss()}. Valorile sunt scrise în tabela
\texttt{benchmarks} după inserarea traiectoriilor, în aceeași tranzacție logică
în care rularea este marcată ca \texttt{completed}. Dacă integrarea numerică sau
persistarea eșuează, runner-ul marchează rularea ca \texttt{failed} și păstrează
mesajul de eroare, astfel încât interfața web să poată diferenția între o
simulare aflată încă în lucru și una care a eșuat definitiv.
